\chapter{Introduction}
\label{ch:intro}

\section{Motivation}
\label{sec:intro-motivation}

Software that performs mathematics is trusted through a fragile stack: tests
that sample a tiny region of the input space, code review that depends on
human attention, and operational monitoring that only observes failures after
they ship. For symbolic computation --- where the artifact is not merely a
number but a transformation of meaning --- this fragility is acute. A single
algebraic mis-step, invisible to unit tests, can propagate through an entire
derivation.

\textsc{mathlib5} takes a different stance. It treats \emph{verification} and
\emph{provenance} as primary engineering outputs, produced alongside every
executable artifact rather than bolted on after the fact. The guiding
principle is stated plainly in the abstract of this work:

\begin{quote}
\itshape Confidence increases when transformations between symbolic notation
and executable artifacts remain inspectable.
\end{quote}

Inspectability has three dimensions in this architecture:
\begin{enumerate}[label=(\roman*)]
  \item \textbf{Syntactic inspectability} --- every stage consumes and emits a
        human-readable artifact (an \textsc{apl} expression, an S-expression,
        a Lean goal, a C{--} term, a receipt).
  \item \textbf{Logical inspectability} --- proof obligations are first-class;
        a result is not emitted unless a prover has discharged its
        obligations (or the gap is explicitly recorded as a \texttt{sorry}).
  \item \textbf{Historical inspectability} --- each transformation writes a
        cryptographically sealed receipt into an append-only WORM ledger, so
        the full lineage of any result is reconstructable and tamper-evident.
\end{enumerate}

\section{The Verified Symbolic Compute Pipeline}
\label{sec:intro-vscp}

The architecture is a pipeline of ten \emph{safety boundaries}. A safety
boundary is a point at which one representation is transformed into the next
under an explicit, checkable contract. The contract may be a parser grammar,
a refinement type, a proof obligation, or a signed receipt. Because each
boundary is independently auditable, the failure of any one layer does not
silently corrupt the rest; it is localized, reported, and sealed.

\begin{table}[H]
\centering
\caption{The ten layers of the Verified Symbolic Compute Pipeline.}
\label{tab:layers}
\begin{tabularx}{\textwidth}{clX}
\toprule
\textbf{\#} & \textbf{Name} & \textbf{Contract / Output} \\
\midrule
1  & Symbolic APL frontend      & Parse \textsc{apl} $\ensuremath{rightarrow}$ typed AST \\
2  & Canonical symbolic IR      & Normalize AST $\ensuremath{rightarrow}$ canonical S-expr \\
3  & Liquid Haskell refinement  & Refinement types discharge structural constraints \\
4  & Lean~4 proof obligations   & Theorems / \texttt{sorry} ledger \\
5  & Verified C{--} execution   & Certified numeric kernel \\
6  & Logic \& static analysis   & FOL resolution + CodeQL checks \\
7  & Policy reasoning (ASP)     & Stable-model policy decisions \\
8  & Optimization \& P/NP swarm & Witness search, sealed skills \\
9  & WORM provenance receipts   & Hash-linked, signed lineage \\
10 & Proof completeness         & Automated \texttt{sorry}-closing \\
\bottomrule
\end{tabularx}
\end{table}

The repository \texttt{SNAPKITTYWEST/mathlib5} is the reference implementation
of this pipeline. Its top-level layout mirrors the layer taxonomy:

\begin{verbatim}
mathlib5/
  kernel/            Layer 5  verified_kernel.cm  (C--)
  layers/apl/        Layer 1  Parser.hs
  layers/sexpr/      Layer 2  AST.hs
  layers/liquid/     Layer 3  Bridge.hs
  layers/hol/lean/   Layer 4  Mathlib5/Core.lean
  layers/fol/        Layer 6  resolution checker
  layers/codeql/     Layer 6  static analysis
  layers/asp/        Layer 7  stable models
  layers/prism-skills/  Layer 8  canonical JSON + WORM
  layers/pnp-attack/ Layer 8  P/NP swarm
  mathlib5-ffi-bridge/   C<->Lean FFI + TheoremCalc
  malice_layer2/    adversarial refutation
  sorrydb/          sorry ledger
  theorems/         closed-form theorems
\end{verbatim}

\section{Contributions}
\label{sec:intro-contributions}

This monograph makes the following contributions:
\begin{enumerate}[label=(\arabic*)]
  \item A coherent ten-layer \emph{safety-boundary} model for symbolic
        computation, with each layer's contract made explicit.
  \item A concrete realization of that model in a mixed-language system
        (\textsc{apl}, Haskell, Lean~4, C{--}, Shell, Nix) that is
        publicly available and reproducible.
  \item A verified numeric kernel expressed in C{--} whose primitives
        (polynomial addition, matrix multiplication, symbolic
        differentiation, algebraic decision, normalization) carry proof
        certificates.
  \item A foreign-function bridge that \emph{cross-validates} the C{--}
        kernel against Lean~4, so two independent paths must agree before a
        result is accepted.
  \item A provenance layer in which every artifact is hash-linked,
        Ed25519-signed, and anchored to an immutable WORM ledger, with this
        very manuscript sealed by the same mechanism.
\end{enumerate}

\section{Organization of this Monograph}
\label{sec:intro-org}

\Cref{part:foundations} (Chapters~\ref{ch:related}--\ref{ch:coherence})
establishes motivation, related work, mathematical preliminaries, and the
coherence/governance model. \Cref{part:pipeline} (Chapters~\ref{ch:vscp}--
\ref{ch:completeness}) walks the ten layers in order, each grounded in the
reference sources. \Cref{part:kernel} (Chapters~\ref{ch:kernel}--
\ref{ch:build}) details the kernel, the FFI bridge, and the build/CI
verification gates. \Cref{part:eco} (Chapters~\ref{ch:security}--
\ref{ch:constellation}) covers the security model and the wider repository
constellation. \Cref{part:eval} (Chapters~\ref{ch:eval}--\ref{ch:conclusion})
presents evaluation, a case study, limitations, and conclusions. Five
appendices reproduce the canonical sources, the metadata schema, a glossary,
and the bibliography.

\section{Reading Guidance}
\label{sec:intro-reading}

A practitioner who wants only the architecture may read
Chapters~\ref{ch:vscp} and \ref{ch:kernel}. A formal-methods reader should
attend to Chapters~\ref{ch:liquid}--\ref{ch:lean} and \ref{ch:completeness}.
A security reader should focus on Chapters~\ref{ch:worm}, \ref{ch:security},
and \ref{ch:constellation}. All readers are encouraged to clone the
repository and follow along with the appendices.
